%%%%%%%%%%%%%%%%%%%%%%%%%%%%%%%%%%%%%%%%%%%%%%%%%%%
%%% DISCUSION
%%%%%%%%%%%%%%%%%%%%%%%%%%%%%%%%%%%%%%%%%%%%%%%%%%%

\chapter{Discusión}
\fancyhead[RE]{\textsc{CAPÍTULO} \thechapter. Discusión}
\label{ch:Discusion}

\noindent Este capítulo analiza y discute en profundidad los resultados obtenidos en la evaluación presentada en el capítulo anterior. 

Voy a realizar un análisis critico, tanto de analisis de resultados como de analisis de errores y decisiones tomadas.

A pesar de las restricciones computacionales y de tiempo, se ha logrado desarrollar un sistema competitivo para la \gls{Tarea}, alcanzando una posición destacada en la clasificación general. Esto demuestra la eficacia de la arquitectura de sistema propuesto, que combina múltiples modelos y técnicas de fusión y re-ranking para maximizar el rendimiento.

Sin embargo, es importante reconocer las limitaciones del enfoque adoptado. La dependencia de modelos preentrenados puede introducir sesgos inherentes a los datos de entrenamiento originales, lo que podría afectar la generalización del sistema a nuevos dominios o lenguajes. Además, la complejidad computacional de los modelos utilizados limita la escalabilidad del sistema para aplicaciones en tiempo real o con grandes volúmenes de datos.

En cuanto a la seleccion de fuentes para los modelos de recuperación, a pesar de intentar abarcar un rango amplio de modelos entrenados en muchos campos distintos, hay un sesgo hacia el MTEB leaderboard, que a pesar de ser un referente, no deja de ser una fuente sesgada y limitada de informacion. Otro punto complejo ha sido la busqueda de modelos preentrenados exclusivamente en español, ya que la mayoría de modelos disponibles están entrenados en inglés o en múltiples idiomas, lo que puede afectar el rendimiento en tareas específicas del español. No se ha encontrado ningun modelo contrastivo entrenado exclusivamente en español, lo que ha obligado a utilizar modelos multilingües que pueden no estar optimizados para las particularidades del idioma.

En cuanto a la figura \ref{fig:zero_shot_map} se observa una clara predominancia de los modelos multilingues\footnote{Cabe destacar qeu la informacion de cuantos lenguajes se habia entrenado viene dada por el ultimo entrenamiento realizado, pero muchos de ellos son ajustes finos de modelos mas generales, por lo general entrenados para multilingue, reduciendo aun más la fiabilidad de las etiquetas monolingues que se muestran en la grafica.}. Los monolingues quedan en un claro segundo puesto, pudiendo estar motivado por la falta de generalización de los modelos exclusivamente monolingues, o quizas por la presencia de anglicismos en multiples terminos de titulos de trabajos presentes hoy en dia en el idioma español.

Continuando con la figurra, se puede observar tambien dos peculiaridades en su distribución. La primera es que hay muchos que estan alineados en vertical, lo que indica que muchos modelos comparten la misma arquitectura base, y por tanto, es probable que compartan tambien limitaciones y sesgos similares (luego comentaremos porque para los tres modelos elegidos para la fusion, hay dos que tienen la misma arquitectura base, e igual no fue una buena decision). La segunda es que, tal y como reportaron la organizacion de la tarea en \cite{talentclef2025_overview}, los modelos parecen presentar un punto dulce en cuanto a tamaño, a partir del cual modelos más grandes presentan empeoramientos en el rendimiento de la tarea. Este punto se puede ver empiricamente entre 500M y 1000M de parametros, donde se observa un pico en el rendimiento, y a partir de ahi, los modelos más grandes tienden a tener un rendimiento inferior. Esto puede deberse a varios factores, discusion en la que no voy a entrar.

Luego se eligieron los 10 mejores para realizar un ajuste fino. Esta claro que la estrategia de entrenamiento muestra que no se podia ajustar a muchas epocas debido al sobreajuste, pero quizas ya con una epoca, que es lo que gemos entrenado, ya se ha sobreajustado (a pesar de que hayan subido las metricas de validación). Además, haber elegido solo los 10 mejores modelos puede haber limitado la diversidad de enfoques y arquitecturas, lo que podría haber afectado la capacidad del sistema para capturar diferentes aspectos de la tarea. Un ejemplo es el modelo \textit{Alibaba-NLP/gte-multilingual-base}, que aunque se encontraba en el puesto 9 en la evaluación zero-shot, tras el ajuste fino no logró mejorar su rendimiento, posiblemente debido a su arquitectura o a la naturaleza de los datos de entrenamiento utilizados. esto lo coloco en el tercer puesto, haciendo que formara parte del sistema final. Quizas se podria haber ampliado la lista a la que se aplico el ajuste fino, para asegurar que modelos que a priori no estaban en el top 10, pero que podrian haber mejorado significativamente tras el ajuste fino, no fueran descartados prematuramente. Relacionado con el ajuste fino, la mejora no ha sido muy sustancial, lo que sugiere que los modelos preentrenados ya estaban bastante bien adaptados a la tarea, o que los datos de entrenamiento disponibles no eran suficientes para lograr mejoras significativas. 

En relacion con los datos de entranmiento, quizas hubiera sido interesante haber investigado acerca de fuentes de datos con los que alimentar el conjunto de datos de entrenamiento, aportando mayor diversidad y cantidad de ejemplos. Esto podría haber ayudado a los modelos a generalizar mejor y a capturar una gama más amplia de variaciones en los títulos de trabajo.

En cuanto a la seleccion de modelos para la fusion, se seleccionaron los tres mejores, pero dos de ellos compartian la misma arquitectura base (\textit{??}). Esto podría haber limitado la diversidad de enfoques y perspectivas en la fusión, ya que los modelos con arquitecturas similares tienden a cometer errores similares. Una estrategia alternativa podría haber sido seleccionar modelos con arquitecturas más diversas, incluso si su rendimiento individual era ligeramente inferior, para maximizar la complementariedad entre ellos.

Otro sesgo que tambien tenemos en la fusion es que se han seleccionado solo tres modelos, y a partir de ahi se ha realizado la fusion. Quizas hubiera sido interesante haber explorado combinaciones de mas modelos, o haber probado varias combinaciones de los modelos ya ajustados, para ver si alguna combinacion alternativa hubiera podido mejorar el rendimiento final. En cualquier caso, la mejora de la fusion no ha llegado ni a un punto porcentual, mostrando que quizas la estrategia de fusion no ha sido la mas adecuada, o que los modelos seleccionados no eran lo suficientemente complementarios entre si.

Para la fase de reranking se ha intentado proceder como con los modelos de retrieval: hacer una evaluacion masiva de modelos para poder seleccionar con mejor criterio el modelo cross enconder. Sin embargo, en este caso no contamos con un ranking publico de mejores modelos para la tarea (mteb tiene fundamentalmente bi encoders y modelos densos). Ya al hacer el retrieval me di cuenta de que en hugging face no todos los modelos son etiquetados de forma correcta (bi encoders, cross enconders, o modelos de retrieval, modelos de ranking), produciendo una vision sesgada en la busqueda de modelos, ya que solo voy a tener acceso a los que si hayan etiquetado correctamente su modelo. es por ello que el numero de modelos conseguidos, 90, no es ni de lejos representativo con todos los que deberian aparecer ahi fuera. 

Tambien es cierto que el haber decidido trabajar con sentence transformers ha limitado en gran medida el acceso a los modelos disponibles.

He tratado en el reranking de obtener primero una vision general del rendimiento y adecuacion de los modelos a la tarea, viendo su comportamiento en la reordenacion de distintas cantidades 5,10,25 50 y 100. Esto ha mostrado que los reordenamientos de valores grandes empeoraban los resultados ya conseguidos con el ajuste fino y la fusion. por ello, al ver que la mayoria de los modelos mejoraba la metrica de la fusion con una reordenacion de un top pequeño, se decidió hacer un analisis mas detallado de cómo de acusada era la degradación. En la figura \ref{fig:top10_reranking_individual} se ve cómo parece que para la tarea en cuestion, los modelos en general tienden a empeorar la metrica cuando k es mayor o igual a 20. 

Sorprende aqui mucho el comportamiento de modelos como posih-reranker-base-mse, que a pesar de presentar la peor metrica para k=5, es el que menos degradacion sufre cuando k aumenta. Un comportamiento similar sufre mmarco-mMiniLMv2-L12-H384-v1, que en este caso incluso es el unico que mejora su rendimiento para k=10, aunque luego se degrade igualmente. Este último es el que resultados mas estables reporta, siendo una buena opcion a considerar si se buscara generalización. Debido a que solo se buscaba la maximización de la metrica MAP, es por ello que se ha seleccionado el mejor simplemente, aunque es probable que esta decision haya repercutido en la metrica sobre el conjunto de test.